\begin{table}
\caption{Tabela 2. Correlação parcial (controles: latitude, água e altitude) entre a fração protegida (TI ou UC federal) da célula e cada resposta, com IC 95% por blocos de 1°.}
\label{tab:11_territorio::associacao}
\begin{tabular}{llrl}
\toprule
Bacia & Resposta & Células & ρ parcial [IC 95\%] \\
\midrule
Xingu & perda de floresta 2000–2024 & 668 & -0,66 [-0,74; -0,55] \\
Xingu & densidade de focos 2003–2024 & 668 & -0,68 [-0,75; -0,57] \\
Xingu & LST diurno, estação seca (MODIS) & 668 & -0,68 [-0,76; -0,55] \\
Xingu & aquecimento do LST entre as décadas (MODIS) & 668 & -0,53 [-0,64; -0,41] \\
Xingu & máxima do ar, estação seca (ERA5) & 668 & -0,07 [-0,23; +0,07] \\
Xingu & entropia da chuva (CHIRPS) & 668 & +0,31 [+0,14; +0,47] \\
Tapajós & perda de floresta 2000–2024 & 652 & -0,51 [-0,61; -0,38] \\
Tapajós & densidade de focos 2003–2024 & 652 & -0,39 [-0,53; -0,25] \\
Tapajós & LST diurno, estação seca (MODIS) & 652 & -0,51 [-0,64; -0,34] \\
Tapajós & aquecimento do LST entre as décadas (MODIS) & 652 & -0,50 [-0,60; -0,36] \\
Tapajós & máxima do ar, estação seca (ERA5) & 652 & -0,08 [-0,30; +0,18] \\
Tapajós & entropia da chuva (CHIRPS) & 652 & +0,07 [-0,12; +0,31] \\
\bottomrule
\end{tabular}
\end{table}
